\notechapter{N-20210204}{Sets, Classes, and the First Warnings About Size}

\section{Why classes enter before functions}

The basic language of sets, classes, indexed families, and functions begins with a distinction that is easy to underestimate: sets may be elements of other collections, whereas proper classes cannot be treated in the same way without recreating the classical paradoxes.

\begin{definition}
A \vocab{set} is a class which can be an element of another class.  A class which is not a set is called a \vocab{proper class}.
\end{definition}

This is not meant to be the whole foundation of set theory, but it explains the working convention of the note: some collections are small enough to be members of other collections, while others are too large and must remain proper classes.

\begin{caution}
The distinction between sets and proper classes is formally essential: it is one of the mechanisms that prevents the usual paradoxes.  In a class theory such as NBG, sets are precisely the classes that may occur as elements of classes; proper classes may not be elements.
\end{caution}

\section{Russell's class as the first warning}

The first serious example is Russell's construction.  Consider
\[
  M = \{X \mid X \text{ is a set and } X \notin X\}.
\]
At first sight this is a natural class: many sets are not members of themselves.  For instance, the set of all books is not itself a book.

\begin{proposition}
The class \(M\) is not a set.
\end{proposition}

\begin{proof}
Assume, for contradiction, that \(M\) is a set.  Then it makes sense to ask whether \(M \in M\).  By the defining condition of \(M\),
\[
  M \in M \quad \Longleftrightarrow \quad M \notin M,
\]
which is impossible.  Hence \(M\) must be a proper class.
\end{proof}

This is the first place where the note is not merely listing definitions.  The example shows why unrestricted comprehension cannot simply produce a set from every condition.

\section{Power sets and indexed families}

\begin{definition}
The \vocab{power set axiom} asserts that for every set \(A\), the class
\[
  \mathcal P(A)=\{X \mid X\subseteq A\}
\]
of all subsets of \(A\) is itself a set.  It is also common to write \(2^A\) for this power set.
\end{definition}

The next construction is an indexed family.  Given a family of classes or sets
\[
  \{A_i \mid i\in I\},
\]
its union and intersection are written
\[
  \bigcup_{i\in I} A_i = \{x \mid x\in A_i \text{ for some } i\in I\},
\]
and
\[
  \bigcap_{i\in I} A_i = \{x \mid x\in A_i \text{ for all } i\in I\}.
\]

\begin{remark}
The indexing collection \(I\) need not be finite or countable.  This is a useful mental shift: the notation \(\{A_i\}_{i\in I}\) is not just a long list.  It may be a genuinely infinite family.
\end{remark}

\begin{caution}
If \(I\) is allowed to be a proper class, then extra care is needed.  In ordinary ZFC, one usually speaks of set-indexed families.  In a class-theoretic setting, class-indexed families can be discussed, but one must check separately when their unions or intersections are sets.
\end{caution}

If \(A\) and \(B\) are classes, the relative complement of \(A\) in \(B\) is
\[
  B\setminus A = \{x\mid x\in B \text{ and } x\notin A\}.
\]
If all classes under discussion are subsets of a fixed universe \(U\), then \(U\setminus A\) is often denoted \(A^c\).

\section{Set identities}

The standard set identities take the following form:
\begin{proposition}
Let \(A\), \(B\), and \(B_i\) be subsets of a fixed universe \(U\).  Then
\[
 A\cap\left(\bigcup_{i\in I} B_i\right)
 = \bigcup_{i\in I}(A\cap B_i),
\]
\[
 A\cup\left(\bigcap_{i\in I} B_i\right)
 = \bigcap_{i\in I}(A\cup B_i),
\]
\[
 \left(\bigcup_{i\in I} A_i\right)^c
 = \bigcap_{i\in I} A_i^c,
 \qquad
 \left(\bigcap_{i\in I} A_i\right)^c
 = \bigcup_{i\in I} A_i^c.
\]
Moreover, \(A\cup B=B\cup A\), \(A\subseteq B\) means every element of \(A\) lies in \(B\), and \(A\cap B=A\) when \(A\subseteq B\).
\end{proposition}

\begin{proof}
All of these follow by chasing membership.  For instance, an element \(x\) lies in \(A\cap(\bigcup_i B_i)\) exactly when \(x\in A\) and \(x\in B_i\) for at least one \(i\), which is exactly the condition that \(x\in \bigcup_i(A\cap B_i)\).  The De Morgan laws are proved similarly after replacing ``there exists'' by ``for all'' under negation.
\end{proof}

\section{Images, inverse images, and functions}

Let \(f:A\to B\) be a function.  If \(S\subseteq A\), the image of \(S\) under \(f\) is
\[
  f(S)=\{b\in B\mid b=f(a)\text{ for some }a\in S\}.
\]
If \(T\subseteq B\), the inverse image of \(T\) is
\[
  f^{-1}(T)=\{a\in A\mid f(a)\in T\}.
\]

\begin{example}
For \(f:\mathbb R\to\mathbb R\), \(f(x)=x^2\), we have
\[
  f(\mathbb R)=[0,\infty),
  \qquad
  f^{-1}(\{9\})=\{-3,3\}.
\]
\end{example}

A map \(f:A\to B\) is \vocab{injective} if \(f(a)=f(a')\) implies \(a=a'\).  It is \vocab{surjective} if \(f(A)=B\).  It is \vocab{bijective} if it is both injective and surjective.

\begin{proposition}
For composable maps \(A\xrightarrow{f}B\xrightarrow{g}C\):
\begin{enumerate}[(i)]
  \item if \(f\) and \(g\) are injective, then \(g\circ f\) is injective;
  \item if \(f\) and \(g\) are surjective, then \(g\circ f\) is surjective;
  \item if \(g\circ f\) is injective, then \(f\) is injective;
  \item if \(g\circ f\) is surjective, then \(g\) is surjective.
\end{enumerate}
\end{proposition}

\section{One-sided inverses}

One-sided inverses recognize injective and surjective maps.

\begin{theorem}
Let \(f:A\to B\) be a function and assume \(A\neq\varnothing\).
\begin{enumerate}[(i)]
  \item \(f\) is injective if and only if there is a map \(g:B\to A\) such that \(g\circ f=1_A\).
  \item \(f\) is surjective if and only if there is a map \(h:B\to A\) such that \(f\circ h=1_B\), provided the required choices of elements from the fibres of \(f\) can be made.
\end{enumerate}
\end{theorem}

\begin{proof}
For (i), suppose \(f\) is injective.  For \(b\in f(A)\), there is a unique \(a\in A\) with \(f(a)=b\).  Choose once and for all an element \(a_0\in A\).  Define
\[
  g(b)=
  \begin{cases}
    a, & b\in f(A)\text{ and }f(a)=b,\\
    a_0, & b\notin f(A).
  \end{cases}
\]
Then \(g(f(a))=a\) for all \(a\in A\).  Conversely, if \(g\circ f=1_A\), then \(f(a)=f(a')\) implies
\[
  a=g(f(a))=g(f(a'))=a',
\]
so \(f\) is injective.

For (ii), if \(h\) exists and \(f\circ h=1_B\), then every \(b\in B\) equals \(f(h(b))\), so \(f\) is surjective.  Conversely, if \(f\) is surjective, each fibre \(f^{-1}(b)\) is nonempty.  Choosing \(h(b)\in f^{-1}(b)\) gives a right inverse.
\end{proof}

\begin{caution}
The last choice of one element from each fibre is exactly where the Axiom of Choice may enter.  This point will reappear in the chapter on choice and cardinal numbers.
\end{caution}

\section{Set-level bookkeeping behind functions}

The calculations in this note are elementary only after the ambient set has been fixed.  Russell's class shows why the notation
\[
  \{x\mid P(x)\}
\]
cannot be used without asking whether the collection is actually a set.  Once a set \(A\) is available, the usual constructions are safe: subsets live in \(\mathcal P(A)\), indexed unions and intersections are formed from a family \((A_i)_{i\in I}\), and membership proofs reduce identities such as De Morgan's laws to quantifier manipulations.

For functions the same bookkeeping reappears in a more operational form.  The inverse image \(f^{-1}(T)\) is always defined for \(T\subseteq B\), while the image \(f(S)\) depends on the part of the domain actually hit.  Injectivity is detected by a left inverse \(g\circ f=1_A\); surjectivity is detected by a right inverse \(f\circ h=1_B\).  The construction of \(g\) for an injection uses only a single default element of \(A\), whereas the construction of \(h\) for a surjection asks for one chosen element from each fibre \(f^{-1}(b)\).  That fibre-by-fibre choice is the exact point to remember when the later notes discuss the Axiom of Choice and cardinal comparison.
